\section{拉普拉斯变换}\label{sec:Laplace_Transform}

十八世纪，欧拉在研究微分方程时，首先使用了积分变换的方法，将微分方程转化为代数方程来求解，随后，拉格朗日和拉普拉斯也使用了类似的方法。其中，拉普拉斯所使用的积分变换为$y_s=\int \mathrm{e}^{-sx}\phi(x)\,dx$，与现代的拉普拉斯变换最为相似，但当时并未以拉普拉斯变换命名。十九世纪，Petzval在其著作中首次系统地考察了这一方法，但他的学生Spitzer强烈主张将这一变换命名为拉普拉斯变换，这影响到几位很有影响力的数学家：Boole，Poincaré和Doetsch，他们在著作中均将这种方法称为“拉普拉斯变换”，使得这个名称被逐渐巩固。

%二十世纪初，Heaviside在研究电磁学时提出了一种符号运算的方法，将微分方程中的微分算子$\dfrac{d}{dt}$替换为符号$p$（或$D$），从而将微分方程转化为代数方程来求解。这种方法直观、快捷但缺乏严格的数学基础，随着数学家对拉普拉斯变换的研究不断深入，到1950年，拉普拉斯变换已完全取代Heaviside演算。

鉴于拉普拉斯变换的应用场景中使用角频率明显多于使用频率，我们在本书中采用角频率$\omega$作为拉普拉斯变换的变量；另外，本章默认$u(0)=1$。

\subsection{拉普拉斯变换的引出}\label{sub:introduction_of_laplace}
在\ref{sub:introduction_of_fourier}中，我们介绍了傅里叶变换：\begin{align*}
    \mathcal{F} f(\omega)=\int_{-\infty}^{\infty}f(t)\mathrm{e}^{-\mathrm{i}\omega t}\,dt
\end{align*}

在不涉及收敛性的问题时，可以将$\omega$推广为复数$s=\sigma+\mathrm{i}\omega$，根据含参变量积分的求导法则，这样的积分变换可以对$s$求导，从而是全纯函数。这个积分变换就是\textbf{双边拉普拉斯变换}：
\begin{definition}[双边拉普拉斯变换]
    设函数$f$局部可积，$\exists a,b\in\mathbb{R},a<b$（我们允许$a=-\infty$和$b=\infty$），使得$f(t)\mathrm{e}^{-at}$和$f(t)\mathrm{e}^{-bt}$均绝对可积，则对$\text{Re}\ s\in[a,b]$，$f$的双边拉普拉斯变换定义为
    \[\mathcal{L}f(s)=\int_{-\infty}^{\infty}f(t)\mathrm{e}^{-st}\,dt\]
\end{definition}
“局部可积”即$f$在任意有限区间上可积，这个条件是为了保证在添加适当的收敛因子$\mathrm{e}^{-st}$后，拉普拉斯积分收敛，从而拉普拉斯变换有定义。

在$\sigma>0$时，这个收敛因子会使得$t\to -\infty$时$|\mathrm{e}^{-st}|=\mathrm{e}^{-\sigma t}$快速增长，从而对所变换的函数提出很高的要。实际上，更多时候我们仅使用\textbf{单边拉普拉斯变换}，定义如下\cite{folland}：
\begin{definition}[单边拉普拉斯变换]
    记\[\mathcal{E}=\left\{f:[0,\infty)\to\mathbb{C}\left|f\in PC[0,\infty),\exists C>0,a\in\mathbb{R},|f(t)|<Ce^{at},\forall t\geq 0\right.\right\}\]
    其中$PC[0,\infty)$表示$[0,\infty)$上的分段连续函数。
    对$f\in\mathcal{E}$，其拉普拉斯变换定义为
    \[\mathcal{L} f(s)=\int_{0}^{\infty}f(t)\mathrm{e}^{-st}\,dt,\quad s=\sigma+\mathrm{i}\omega\]
    其中$s$的实部$\sigma$必须满足$\sigma>a$，将它称为拉普拉斯变换的\textbf{收敛域}，直线$\sigma=a$称为\textbf{收敛轴}。\\
    如果$f(t)$是时域信号，$\mathcal{L} f(s)$称为\textbf{复频域信号}。
    我们仍记$\mathcal{L} f(s)$为$F(s)$或$\hat{f}(s)$（在涉及傅里叶变换时，我们会避免使用它们或加以说明），并记$f\overset{\mathcal{L} }{\longleftrightarrow} F(s)$。
\end{definition}
在$\mathcal{E} $的定义中我们对$f$的定义域做了要求，但实际上只要$f$是因果信号（见\ref{sub:term_in_ODE}，$f(t)=f(t)u(t)$），就可以认为$f\in\mathcal{E}$，后文中不会特意区分$f$的定义域。

在$\mathcal{E} $上定义拉普拉斯变换，是因为对于$f\in\mathcal{E} $，当$\sigma>a$时，给出拉普拉斯变换的积分是绝对收敛的：
\begin{align*}
    \int_0^{\infty}|f(t)\mathrm{e}^{-st}|\,dt\leq \int_0^{\infty}Ce^{at}\mathrm{e}^{-\sigma t}\,dt=\frac{C}{\sigma-a}<\infty
\end{align*}
因此在导数存在的情况下可以交换积分与求导的次序，从而$\mathcal{L} f(s)$是$\sigma>a$上的全纯函数：
\begin{align*}
    \frac{d}{ds}\mathcal{L} f(s)&=\frac{d}{ds}\int_0^{\infty}f(t)\mathrm{e}^{-st}\,dt\\
    &=\int_0^{\infty}f(t)\frac{\partial}{\partial s}(\mathrm{e}^{-st})\,dt\\
    &=-\int_0^{\infty}tf(t)\mathrm{e}^{-st}\,dt\\
    &=-\mathcal{L} [tf(t)](s)
\end{align*}
这正是拉普拉斯变换的复频域微分性质。

根据全纯函数的唯一延拓性，如果$\mathcal{L} f(s)$可以延拓至$\sigma\leq a$的一些区域，则延拓是唯一的，尽管这时积分不一定存在，我们也认为延拓所得的定义域更大的函数是$f$的拉普拉斯变换。
\begin{example}[多项式的拉普拉斯变换]
    设$f(t)=t^\nu(\nu\geq 0)$，则$f\in\mathcal{E} $，且
    \begin{align*}
        \mathcal{L} f(s)&=\int_0^{\infty}t^\nu \mathrm{e}^{-st}\,dt\\
        &=\frac{1}{s^{\nu+1}}\int_0^{\infty}x^\nu \mathrm{e}^{-x}\,dx&(x=st)\\
        &=\frac{\Gamma(\nu+1)}{s^{\nu+1}},\quad \sigma>0
    \end{align*}
尽管积分只在$\sigma>0$时收敛，但$\mathcal{L} f(s)$可以延拓为$\mathbb{C}\backslash\{0\}$上的全纯函数。
\begin{remark}
    其实，这个公式对实部大于$-1$的$\nu$都成立。尽管$-1<\text{Re}\ \nu<0$时，$f(t)=t^\nu$\\
    在$t=0$处无界，但由于
\[\left|\int_{0}^{1}t^{\nu}\mathrm{e}^{-st}\,dt\right|\leq \int_0^1 t^\nu\,dt=\frac{1}{\nu+1},\sigma>0\]
\[\left|\int_{1}^{\infty}t^{\nu}\mathrm{e}^{-st}\,dt\right|\leq \int_1^\infty \mathrm{e}^{-\sigma t}\,dt=\frac{\mathrm{e}^{-\sigma}}{\sigma},\sigma>0\]
其拉普拉斯变换是存在的。因此，我们可以认为$t^\nu\in\mathcal{E},\text{Re}\ \nu>-1$.
\end{remark}
\end{example}

在$\mathcal{E} $中，卷积总是良定义的：
\begin{proposition}[$\mathcal{E} $中的卷积]
    设$f,g\in\mathcal{E} $，则它们的卷积为
    \[(f*g)(t)=\begin{cases}
    0                   & \text{if } t<0 \\
    \int_0^t f(x)g(t-x)\,dx & \text{if } t\geq 0
    \end{cases}\]
\end{proposition}
这是因为$x\geq 0,t-x\geq 0$时，$f(x),g(t-x)$才不恒为0，这相当于积分区间为$[0,t]$。

\subsection{拉普拉斯变换的运算性质}\label{sub:laplace_operation}

\begin{proposition}[拉普拉斯变换的性质]
    设$f,g\in\mathcal{E} $，$a,b\in\mathbb{C}$，则：\begin{itemize}
        \item 线性：在$f,g$的公共收敛域内，有\[\mathcal{L} [af(t)+bg(t)](s)=a\mathcal{L} f(s)+b\mathcal{L} g(s)\]
        \item 时移：\[\mathcal{L} [f(t-t_0)u(t-t_0)](s)=\mathrm{e}^{-s t_0}\mathcal{L} f(s)\]
        \item 频移：\[\mathcal{L} [\mathrm{e}^{a t}f(t)](s)=\mathcal{L} f(s-a)\]
        \item 伸缩：设$a>0$，则\[\mathcal{L} [f(at)](s)=\frac{1}{a}\mathcal{L} f\left(\frac{s}{a}\right)\]
        \item 时域微分：设$f\in C^1[0,\infty),f'\in\mathcal{E} $，则\[\mathcal{L} [f'(t)](s)=s\mathcal{L} f(s)-f(0)\]
        $f'(t)$在0处的定义为$f$在0处的右导数：$$f'(0):=\lim_{h\to 0+}\dfrac{f(h)-f(0)}{h}$$
        注意我们必须要求$f$是连续的，否则微积分基本定理需要分段使用，会使公式更加复杂。得到时域微分性质至少应该要求$f$是连续且分段光滑的。
        \item 复频域微分：\[\mathcal{L} [t f(t)](s)=-\frac{d}{ds}\mathcal{L} f(s)\]
        \item 时域积分：\[\mathcal{L} \left[\int_0^t f(x)\,dx\right](s)=\frac{1}{s}\mathcal{L} f(s)\]
        \item 复频域积分：设$\dfrac{f(t)}{t}\in\mathcal{E} $，则\[\mathcal{L} \left[\frac{f(t)}{t}\right](s)=\int_s^{\infty}\mathcal{L}f(x)\,dx\]
        此处积分路径为复平面上满足$Im(s)=\omega$有界的任意从$s$到无穷远处的路径
    \end{itemize}
\end{proposition}
\begin{remark}
    很明显，时域微分性质和复频域微分性质都可以对更高阶的导数进行推广，其中，时域微分性质推广为\begin{align*}
\mathcal{L} [f^{(n)}(t)](s)&=s^n\mathcal{L} f(s)-\sum_{k=0}^{n-1}s^{n-1-k}f^{(k)}(0)\\
&=s^n\mathcal{L} f(s)-s^{n-1}f(0)-s^{n-2}f'(0)-\cdots -f^{(n-1)}(0)
\end{align*}
\end{remark}
\begin{proof}
    1.线性性：直接得自积分的线性性。\\
    2.时移：\begin{align*}
        \mathcal{L} [f(t-t_0)u(t-t_0)](s)&=\int_0^{\infty}f(t-t_0)u(t-t_0)\mathrm{e}^{-st}\,dt\\
        &=\int_{t_0}^{\infty}f(t-t_0)\mathrm{e}^{-st}\,dt\\
        &=\mathrm{e}^{-s t_0}\int_0^{\infty}f(x)\mathrm{e}^{-s x}\,dx&( x=t-t_0)\\
        &=\mathrm{e}^{-s t_0}\mathcal{L} f(s)
    \end{align*}
    3.频移：\begin{align*}
        \mathcal{L} [\mathrm{e}^{a t}f(t)](s)=\int_0^{\infty}\mathrm{e}^{a t}f(t)\mathrm{e}^{-st}\,dt=\int_0^{\infty}f(t)\mathrm{e}^{-(s-a)t}\,dt=\mathcal{L} f(s-a)
    \end{align*}
    4.伸缩：\begin{align*}
        \mathcal{L} [f(at)](s)&=\int_0^{\infty}f(at)\mathrm{e}^{-st}\,dt\\
        &=\int_0^{\infty}f(x)\mathrm{e}^{-\frac{s}{a}x}\frac{1}{a}\,dx&(x=at)\\
        &=\frac{1}{a}\mathcal{L} f\left(\frac{s}{a}\right)
    \end{align*}
    5.时域微分：\begin{align*}
        \mathcal{L} [f'(t)](s)&=\int_0^{\infty}f'(t)\mathrm{e}^{-st}\,dt\\
        &=\evalat{\left(f(t)\mathrm{e}^{-st}\right)}{0}{\infty}+\int_0^{\infty}sf(t)\mathrm{e}^{-st}\,dt\\
        &=s\mathcal{L} f(s)-f(0)
    \end{align*}
    6.复频域微分：已经在本节前面证明过。\\
    7.时域积分：利用时域微分性质，有
    \begin{align*}
        \mathcal{L} \left[\int_0^t f(x)\,dx\right](s)=\frac{1}{s}\mathcal{L} \left[\frac{d}{dt}\left(\int_0^t f(x)\,dx\right)\right](s)=\frac{1}{s}\mathcal{L} f(s)
    \end{align*}
    8.复频域积分：利用复频域的微分性质，有
    \begin{align*}
        \frac{d}{ds}\mathcal{L} \left[\frac{f(t)}{t}\right](s)&=-\mathcal{L} f(s)\\
        \mathcal{L} \left[\frac{f(t)}{t}\right](s)&=\int_s^{\infty}\mathcal{L}f(x)\,dx
    \end{align*}
\end{proof}

\begin{theorem}[卷积定理]
    设$f,g\in\mathcal{E} $，则
    \[\mathcal{L} [(f*g)(t)](s)=\mathcal{L} f(s)\cdot \mathcal{L} g(s)\]
\end{theorem}
\begin{proof}
与前面傅里叶变换的卷积定理证明类似：
    \begin{align*}
        \mathcal{L} [(f*g)(t)](s)&=\int_0^{\infty}(f*g)(t)\mathrm{e}^{-st}\,dt\\
        &=\int_0^{\infty}\left(\int_0^t f(x)g(t-x)\,dx\right)\mathrm{e}^{-st}\,dt\\
        &=\int_0^{\infty}\int_0^{\infty}f(x)g(t-x)u(t-x)\mathrm{e}^{-st}\,dx\,dt\\
        &=\int_0^{\infty}f(x)\left(\int_x^{\infty}g(t-x)\mathrm{e}^{-st}\,dt\right)\,dx\\
        &=\int_0^{\infty}f(x)\mathrm{e}^{-sx}\left(\int_0^{\infty}g(y)\mathrm{e}^{-sy}\,dy\right)\,dx&(y=t-x)\\
        &=\mathcal{L} f(s)\cdot \mathcal{L} g(s)
    \end{align*}
\end{proof}

以上所有性质的证明都与傅里叶变换类似，实际上我们不需要做这种重复性的证明工作。给定$s=\sigma+\mathrm{i}\omega$，令$\tilde{f}(t)=f(t)\mathrm{e}^{-\sigma t}$，则$\mathcal{L} f(s)=\mathcal{F} \tilde{f}(\omega)$，只要$\tilde{f}\in L^1(\mathbb{R} )$，就可以将傅里叶变换的性质直接应用到拉普拉斯变换上。由于我们在$\mathcal{E} $中要求了$f(t)$的分段连续性，在收敛域$\sigma>a$内，$\tilde{f}(t)$总是属于$L^1(\mathbb{R})$的。下面给出一个例子：
%\begin{example}[时移性质]
%    设$f\in\mathcal{E} ,t_0>0$，令$\tilde{f}(t)=f(t)\mathrm{e}^{-\sigma t}$，则
%    \begin{align*}
%        \mathcal{L}[f(t-t_0)u(t-t_0)](s)&=\mathcal{F} \left[\tilde{f}(t-t_0)u(t-t_0)\right](\omega)\\
%        &=\mathrm{e}^{-\mathrm{i}\omega t_0}\mathcal{F} \tilde{f}(\omega)\\
%        &=\mathrm{e}^{-s t_0}\mathcal{L} f(s)
%    \end{align*}
%\end{example}
\begin{example}[卷积定理]
设$f,g\in\mathcal{E} $，令$\tilde{f}(t)=f(t)\mathrm{e}^{-\sigma t}, \tilde{g}(t)=g(t)\mathrm{e}^{-\sigma t}$，注意到\begin{align*}
    (\tilde{f}*\tilde{g})(t)&=\int_{-\infty}^{\infty}\tilde{f}(x)\tilde{g}(t-x)\,dx\\
    &=\int_0^{t}f(x)\mathrm{e}^{-\sigma x}g(t-x)\mathrm{e}^{-\sigma (t-x)}\,dx\\
    &=\mathrm{e}^{-\sigma t}\int_0^{t}f(x)g(t-x)\,dx\\
    &=\mathrm{e}^{-\sigma t}(f*g)(t)
\end{align*}
因此有
    \begin{align*}
        \mathcal{L} \left[f*g\right](s)=\mathcal{F} \left[f*g(t)\mathrm{e}^{-\sigma t}\right](\omega)=\mathcal{F} \left[\tilde{f}*\tilde{g}\right](\omega)=\mathcal{F} \tilde{f}(\omega)\cdot \mathcal{F} \tilde{g}(\omega)=\mathcal{L} f(s)\cdot \mathcal{L} g(s)
    \end{align*}
\end{example}
其余的证明，请读者自行完成。

\subsection{常用函数的拉普拉斯变换}\label{sub:ordinary_laplace}

\begin{example}[指数多项式]
    设$f(t)=t^\nu \mathrm{e}^{a t}u(t),\text{Re}\ \nu>-1$，利用拉普拉斯变换的频移性质，有
    \begin{align*}
        \mathcal{L} f(s)&=\mathcal{L}\left[t^\nu \mathrm{e}^{a t}u(t)\right](s)\\
        &=\mathcal{L} \left[t^\nu u(t)\right](s-a)\\
        &=\frac{\Gamma(\nu+1)}{(s-a)^{\nu+1}},\sigma>a
    \end{align*}
    根据有理函数的理论，任意真分式函数都可以表示为若干个此类函数的线性组合，因此可以说$f$是指数多项式等价于$\mathcal{L} f$是真分式函数
\end{example}

\begin{example}[正弦、余弦函数]
    首先，$\cos(\omega_0 t),\sin(\omega_0 t)\in\mathcal{E}$，因为$$\forall a>0,|\cos(\omega_0 t)|,|\sin(\omega_0 t)|<\mathrm{e}^{a t}$$
    下面给出两种求解方法。\\
    方法一：利用欧拉公式化归为指数函数：
    \begin{align*}
        \mathcal{L}[\cos(\omega_0 t)](s)&=\mathcal{L}[\cos(\omega_0 t)u(t)](s)\\
        &=\frac{1}{2}\mathcal{L} [\mathrm{e}^{\mathrm{i}\omega_0 t}u(t)](s)+\frac{1}{2}\mathcal{L} [\mathrm{e}^{-\mathrm{i}\omega_0 t}u(t)](s)\\
        &=\frac{1}{2}\left(\frac{1}{s-\mathrm{i}\omega_0}+\frac{1}{s+\mathrm{i}\omega_0}\right)\\
        &=\frac{s}{s^2+\omega_0^2},\sigma>0
    \end{align*}
    \begin{align*}
        \mathcal{L}[\sin(\omega_0 t)](s)&=\mathcal{L} [\sin(\omega_0 t)u(t)](s)\\
        &=\frac{1}{2\mathrm{i}}\mathcal{L} [\mathrm{e}^{\mathrm{i}\omega_0 t}u(t)](s)-\frac{1}{2\mathrm{i}}\mathcal{L} [\mathrm{e}^{-\mathrm{i}\omega_0 t}u(t)](s)\\
        &=\frac{1}{2\mathrm{i}}\left(\frac{1}{s-\mathrm{i}\omega_0}-\frac{1}{s+\mathrm{i}\omega_0}\right)\\
        &=\frac{\omega_0}{s^2+\omega_0^2},\sigma>0
    \end{align*}
    方法二：利用微分性质：
    \begin{align*}
        -a^2\mathcal{L}[\cos(\omega_0 t)](s)&=\mathcal{L} [\cos''(\omega_0 t)](s)\\
        &=s^2\mathcal{L} [\cos(\omega_0 t)](s)-s\cos(0)-\cos'(0)\\
        &=s^2\mathcal{L} [\cos(\omega_0 t)](s)-s
    \end{align*}
    因此\[\mathcal{L} [\cos(\omega_0 t)](s)=\frac{s}{s^2+\omega_0^2}\]
    同理，有\[\mathcal{L} [\sin(\omega_0 t)](s)=\frac{\omega_0}{s^2+\omega_0^2}\]
\end{example}

\begin{example}[双曲函数]
    因为
    \[\sinh(\omega_0 t)=\frac{\mathrm{e}^{\omega_0 t}-\mathrm{e}^{-\omega_0 t}}{2},\cosh(\omega_0 t)=\frac{\mathrm{e}^{\omega_0 t}+\mathrm{e}^{-\omega_0 t}}{2}\]
    所以
    \begin{align*}
        \mathcal{L} \left[\sinh(\omega_0 t)\right](s)&=\frac{1}{2}\mathcal{L} \left[\mathrm{e}^{\omega_0 t}\right](s)-\frac{1}{2}\mathcal{L} \left[\mathrm{e}^{-\omega_0 t}\right](s)\\
        &=\frac{1}{2}\left(\frac{1}{s-\omega_0}-\frac{1}{s+\omega_0}\right)\\
        &=\frac{\omega_0}{s^2-\omega_0^2},\sigma>|\omega_0|
    \end{align*}
    \begin{align*}
        \mathcal{L} \left[\cosh(\omega_0 t)\right](s)&=\frac{1}{2}\mathcal{L} \left[\mathrm{e}^{\omega_0 t}\right](s)+\frac{1}{2}\mathcal{L} \left[\mathrm{e}^{-\omega_0 t}\right](s)\\
        &=\frac{1}{2}\left(\frac{1}{s-\omega_0}+\frac{1}{s+\omega_0}\right)\\
        &=\frac{s}{s^2-\omega_0^2},\sigma>|\omega_0|
    \end{align*}
\end{example}

\begin{example}[采样函数]
    首先，因为$\dfrac{\sin(\omega_0 t)}{t}$有界，所以$\dfrac{\sin(\omega_0 t)}{t}\in\mathcal{E}$。

    结合正弦函数的拉普拉斯变换和拉普拉斯变换的复频域积分性质，有
    \begin{align*}
        \mathcal{L} \left[\frac{\sin(\omega_0 t)}{t}\right](s)&=\int_s^{\infty}\mathcal{L}[\sin(\omega_0 t)](x)\,dx\\
        &=\int_s^{\infty}\frac{\omega_0}{x^2+\omega_0^2}\,dx\\
        &=\evalat{\arctan\frac{x}{\omega_0}}{s}{\infty}\\
        &=\frac{\pi}{2}-\arctan\frac{s}{\omega_0}\\
        &=\arctan\left(\frac{\omega_0}{s}\right),\sigma>0
    \end{align*}
\end{example}

\begin{example}[正弦积分函数]
    $$\frac{\sin(\omega_0 t)}{t}\leq Ce^{a t},a>0\implies \text{Si}(\omega_0 t)\leq \frac{C}{a}\mathrm{e}^{a t}\implies \text{Si}(\omega_0 t)\in\mathcal{E}$$
    利用正弦函数的拉普拉斯变换和拉普拉斯变换的时域积分性质，有
    \begin{align*}
        \mathcal{L} [\text{Si}(\omega_0 t)](s)&=\mathcal{L} \left[\int_0^{\omega_0 t} \frac{\sin(x)}{x}\,dx\right](s)\\
        &=\mathcal{L} \left[\int_0^{t} \frac{\sin(\omega_0 x)}{x}\,dx\right](s)\\
        &=\frac{1}{s}\mathcal{L} \left[\frac{\sin(\omega_0 t)}{t}\right](s)\\
        &=\frac{1}{s}\arctan\left(\frac{\omega_0}{s}\right),\sigma>0
    \end{align*}    
\end{example}

\subsection{*拉普拉斯变换的级数求法}\label{sub:laplace_by_series}
我们知道，多项式的拉普拉斯变换都可以通过$$\mathcal{L} [t^n](s)=\dfrac{\Gamma(n+1)}{s^{n+1}}$$来计算，那么对于一些可以展开为幂级数的函数，不难想到逐项求拉普拉斯变换，下面讨论这种操作的适用范围。在此之前，我们先介绍一个引理：
\begin{lemma}
    设$f\in\mathcal{E} ,|f(t)|\leq Ce^{at}$，则\begin{itemize}
        \item $\forall x>a,\mathcal{L} f(x+iy)\to 0,|y|\to\infty$
        \item $\forall y,\mathcal{L} f(x+iy)\to 0,x\to\infty$
    \end{itemize}
\end{lemma}
\begin{proof}
    1.令$g(t)=f(t)\mathrm{e}^{-x t}$，则$|g(t)|\leq Ce^{-(x-a)t}$，因此$g\in L^1(\mathbb{R} )$，根据Riemann-Lebesgue引理，有
    \[\mathcal{L} f(x+iy)=\mathcal{F} g(y)\to 0,|y|\to\infty\]
    2.利用海涅归并原理，任取实部趋于无穷的序列$\{z_n\}$，满足$\mathcal{L} f(z_n)$存在（这要求$\text{Re}\ z>a$），我们要证明序列$\mathcal{L} f(z_n)\to 0$。由控制收敛定理，有
    \begin{align*}
        &\int_{0}^{\infty}|f(t)\mathrm{e}^{-z_n t}|\,dt\leq \int_0^{\infty}Ce^{at}\mathrm{e}^{-Re(z_n) t}\,dt=C\int_0^{\infty}\mathrm{e}^{-(Re(z_n)-a)t}\,dt=\frac{C}{\text{Re}\ z_n-a}<\infty\\
        &\lim_{n\to\infty}\mathcal{L} f(z_n)=\lim_{n\to\infty}\int_0^{\infty}f(t)\mathrm{e}^{-z_n t}\,dt=\int_0^{\infty}\lim_{n\to\infty}f(t)\mathrm{e}^{-z_n t}\,dt=0
    \end{align*}
\end{proof}

\begin{theorem}[拉普拉斯变换的级数求法]
    设$f\in\mathcal{E} $，且在$[0,\infty)$上有幂级数展开
    \[f(t)=\sum_{n=0}^{\infty}a_n t^n\]
    如果存在$\omega_0>0$，使得级数
    \[\sum_{n=0}^{\infty}|a_n|\frac{\Gamma(n+1)}{\omega_0^{n}}\]
    收敛，则在$\sigma>\omega_0$时，级数
    \[\sum_{n=0}^{\infty}a_n \frac{\Gamma(n+1)}{s^{n+1}}\]
    收敛且等于$\mathcal{L} f(s)$，即$f$的拉普拉斯变换可以通过对其幂级数求拉普拉斯变换得到。
\end{theorem}
\begin{proof}
    由条件知级数$\sum_{n=0}^{\infty}a_n \frac{\Gamma(n+1)}{s^{n+1}}$在$\sigma>\omega_0$时绝对收敛、内闭一致收敛，因此只需证明其和函数等于$\mathcal{L} f(s)$。拉普拉斯变换是一种积分变换，我们需要将求和与积分交换，可以考虑使用控制收敛定理。由于
    \begin{align*}
        &|a_n \Gamma(n+1)\omega_0^{-n}|\to 0, n\to\infty\\
        \implies &|a_n|\leq C\frac{\omega_0^{n}}{\Gamma(n+1)}\\
        \implies &|f(t)|\leq \sum_{n=0}^{\infty}|a_n|t^n\leq C\sum_{n=0}^{\infty}\frac{(\omega_0 t)^{n}}{\Gamma(n+1)}=Ce^{\omega_0 t}\\
        \implies &\forall s,\text{Re}\ s=\omega_0+\phi>\omega_0,|f(t)\mathrm{e}^{-s t}|\leq Ce^{-\phi t}
    \end{align*}
    控制收敛定理的条件成立，换序即证。
\end{proof}

一般地，对于幂次不是整数的情况，可以推广以上结论：
\begin{corollary}*
   设$f\in\mathcal{E} $，且在$[0,\infty)$上有广义幂级数展开
   \[f(t)=\sum_{n=0}^{\infty}a_n t^{n+\alpha}\]
   其中$0\leq\alpha<1$，如果存在$\omega_0>0$，使得级数
   \[\sum_{n=0}^{\infty}|a_n|\frac{\Gamma(n+\alpha+1)}{\omega_0^{n}}\]
   收敛，则在$\sigma>\omega_0$时，级数
   \[\sum_{n=0}^{\infty}a_n \frac{\Gamma(n+\alpha+1)}{s^{n+\alpha+1}}\]
   收敛且等于$\mathcal{L} f(s)$，即$f$的拉普拉斯变换可以通过对其广义幂级数求拉普拉斯变换得到。
\end{corollary}
